\scanpage{30}
% archive sheet number: 29
l'identité). Il transforme $0,1,\infty$ en
\begingroup\makeatletter\tagsleft@true\makeatother
\[
\left\{\begin{aligned}
\sigma0&=2\\
\sigma1&=-1\\
\sigma\infty\unsure{.|}&=\tfrac12
\end{aligned}\right.\tag*{7)}
\]
\endgroup
% slanted
\margin{\unsure{Prouver|Préciser} : le normalisateur de $\unsure{\mathfrak{S}|\mathfrak{G}}_3$ dans $\mathrm{Aut}_S P^1_S$ est $\mathfrak{S}_3\times(\deleted{\unsure{\mathfrak{S}|\sigma}}\mathbb{Z}/2)$}

La \unsure{différente|différent} de $\mathfrak{S}_3$ -- ens des pts de ramification -- est formé des ens. des pts où un des cinq \unsure{opérateurs} $\lambda\mu\nu\varphi\psi$ est de l'inertie.

Pour $\lambda,\mu,\nu$ on trouve les sections $(\deleted{\unsure{0|t}}1,\deleted{\unsure{2|1}}-1)$ $(\infty,1/2)$ $(0,2)$, [\unsure{stable|stables} par $\unsure{S|\sigma}$]. Pour $\varphi$ (et $\psi$, c'est pareil), on trouve $1-\frac1t=t$ ou
\begingroup\makeatletter\tagsleft@true\makeatother
\[
t^2-t+1=0\tag*{8)}
\]
\endgroup
i.e. l'équation cyclotomique $\boxed{\Phi_6(t)=0}$. \deleted{\unsure{la|le|6}}

\deleted{\unsure{définit aussi|définit ainsi|définit un|déduit aussi}} C'est aussi l'ens $\overline{S}$ des pts fixes sous $\sigma$.

[Il est non ramifié \inserted{sur $S$} aux pts $s\in S$ tels que $k(s)$ soit de car. $\ne2,3$.]. $\mathbb{Z}_2=\mathfrak{S}_3/(\mathbb{Z}/3)$ y\unsure{:|} opère, de façon \uline{non triviale} aux pts où il est non ramifié sur $S$.

\margin{\unsure{Ceci suppose|bien supposer|en supposant|les pts} qu'il n'y a pas de car. 2 ou 3 dans $S$}D'autre part, le \unsure{schém.|schéma} \inserted{$G$} des automorphismes de $P^1$ qui laissent $\overline{S}$ \unsure{globalement} invariant est une extension de $\mathbb{Z}_2$ [automorphismes de $\overline{S}$] par $(G_m)_S$ \uline{tordu} ---- [par le rev. cyclotomique]

